\chapter{Literature Review}
\label{ch:literature}

\section{Introduction}

This chapter reviews the literature relevant to the design and development of the
EcoTrace: Intelligent Electronic Waste Recovery and Circular Economy Management
System. The review establishes the theoretical and technological foundation of the
study by examining existing knowledge on electronic waste management, circular
economy principles, intelligent waste management, cloud computing, artificial
intelligence, geographic information systems, and modern software development
technologies. Furthermore, the chapter critically reviews existing electronic waste
management systems, identifies their strengths and limitations, and highlights the
research gap that motivated the development of the proposed EcoTrace system.

The literature review is organized into seven major sections. The first section
discusses the conceptual foundations of electronic waste management and related
concepts. The second examines the technologies adopted in the development of
EcoTrace. The third reviews existing electronic waste management systems reported
in literature. This is followed by a comparative analysis of related systems,
identification of the research gap, and a summary of the chapter.

\section{Conceptual Review}

\subsection{Electronic Waste (E-Waste)}

Electronic waste (e-waste), also known as Waste Electrical and Electronic Equipment
(WEEE), refers to discarded electrical and electronic devices that have reached the
end of their useful life, become obsolete, or are no longer required by users
\citep{forti2024}.

The increasing demand for electronic devices, coupled with rapid technological
innovation, shorter product life cycles, and changing consumer preferences, has
contributed significantly to the rapid growth of global electronic waste generation.
According to the Global E-waste Monitor 2024, worldwide electronic waste
generation continues to increase at a rate faster than formal collection and recycling
activities, creating significant environmental, economic, and public health
challenges \citep{forti2024}.

Electronic waste contains both valuable and hazardous materials. Valuable
components include gold, silver, copper, aluminium, platinum, and rare earth
elements that can be recovered through appropriate recycling processes. However,
e-waste also contains hazardous substances such as lead, mercury, cadmium,
chromium, brominated flame retardants, and other toxic chemicals capable of
contaminating soil, water, and air if improperly managed \citep{who2021}.

Proper electronic waste management therefore requires coordinated activities
involving waste generation, collection, transportation, sorting, repair,
refurbishment, recycling, resource recovery, and safe disposal. Modern information
systems increasingly support these activities through cloud computing, artificial
intelligence, geographic information systems, and mobile technologies.

As illustrated in Figure~\ref{fig:ewaste-trends}, global electronic waste generation
has increased steadily over recent years while formal recycling rates remain
comparatively low.

\begin{figure}[H]
  \centering
  \imageplaceholder{Global Electronic Waste Generation and Recycling Trends chart}
  \caption{Global Electronic Waste Generation and Recycling Trends. Source: Adapted from the Global E-waste Monitor 2024 \citep{forti2024}.}
  \label{fig:ewaste-trends}
\end{figure}

\subsection{Circular Economy}

The circular economy is an economic model that aims to eliminate waste and
maximize the continuous use of resources through strategies such as reuse, repair,
refurbishment, remanufacturing, and recycling \citep{geissdoerfer2017,kirchherr2017}.
Unlike the traditional linear economy, which follows the sequence of
``take-make-dispose,'' the circular economy seeks to maintain the value of
products, materials, and resources for as long as possible while minimizing
environmental pollution and resource depletion.

In recent years, the circular economy has become a key strategy for achieving
sustainable development and responsible resource management. Governments,
international organizations, and industries increasingly recognize circular economy
principles as effective approaches for reducing waste generation, conserving
natural resources, lowering carbon emissions, and promoting sustainable
production and consumption \citep{unep2024}.

Electronic waste management is one of the sectors that benefits significantly
from circular economy practices. Instead of disposing of obsolete electronic
devices, valuable components and materials can be recovered through systematic
collection, repair, refurbishment, remanufacturing, and recycling processes. This
not only reduces the volume of waste sent to landfills but also decreases the
demand for virgin raw materials used in manufacturing new electronic products.

The implementation of circular economy principles requires efficient
coordination among multiple stakeholders, including households, businesses,
collectors, logistics providers, technicians, recyclers, government agencies, and
environmental organizations. Digital technologies such as cloud computing,
artificial intelligence, Internet of Things (IoT), blockchain, and geographic
information systems (GIS) have become important enablers of circular economy
initiatives by improving traceability, transparency, and decision-making
throughout the product lifecycle \citep{esmaeilian2020}.

For electronic waste management, the circular economy extends beyond
recycling. It encourages preventive maintenance, repair, refurbishment, reuse, and
responsible redistribution of electronic devices before considering material
recovery or disposal. These practices contribute to resource conservation while
creating new economic opportunities and reducing environmental impacts.

The EcoTrace system was designed based on circular economy principles by
supporting electronic waste registration, collection, transportation, repair,
refurbishment, recycling, and lifecycle tracking within a single integrated platform.
Through these functionalities, the system promotes responsible resource recovery
and supports sustainable electronic waste management practices.

Figure~\ref{fig:circular-framework} illustrates the relationship between product
design, consumption, collection, repair, refurbishment, recycling, and resource
recovery within a circular economy.

\begin{figure}[H]
  \centering
  \imageplaceholder{Circular Economy Framework for Electronic Waste Management}
  \caption{Circular Economy Framework for Electronic Waste Management. Source: Adapted from UNEP (2024).}
  \label{fig:circular-framework}
\end{figure}

\subsection{Reverse Logistics}

Reverse logistics refers to the planning, implementation, and control of the
efficient movement of products, materials, and information from the point of
consumption back to the point of origin or an appropriate recovery facility for the
purpose of reuse, repair, refurbishment, remanufacturing, recycling, or
environmentally responsible disposal \citep{agrawal2015,rogers1999}.

Unlike traditional forward logistics, which focuses on delivering products from
manufacturers to consumers, reverse logistics manages the flow of products after
they have reached the end of their useful life. The primary objective is to recover
economic value from discarded products while minimizing environmental impacts
through sustainable resource management.

Reverse logistics has become an essential component of modern electronic
waste management because discarded electronic devices contain valuable
materials that can be recovered and reintroduced into manufacturing processes.
Effective reverse logistics systems facilitate the collection, transportation,
inspection, sorting, repair, refurbishment, recycling, and final disposal of
electronic waste in a structured and efficient manner \citep{govindan2015}.

The implementation of reverse logistics requires close coordination among
households, businesses, waste collectors, transport operators, repair technicians,
refurbishment centres, recycling companies, and environmental regulatory
agencies. Digital information systems play a critical role in supporting this
coordination by enabling real-time tracking, centralized data management, route
optimization, inventory monitoring, and performance reporting.

Recent advances in cloud computing, mobile technologies, artificial intelligence,
Internet of Things (IoT), and geographic information systems have significantly
improved reverse logistics operations. Organizations increasingly use intelligent
software platforms to optimize collection schedules, monitor transportation
activities, predict waste generation patterns, and improve resource recovery
efficiency \citep{kazancoglu2021}.

Within the EcoTrace system, reverse logistics forms the operational foundation
of the electronic waste management process. The system supports waste
registration, pickup scheduling, collection management, transportation tracking,
repair, refurbishment, recycling, and lifecycle monitoring through an integrated
cloud-based platform. These functionalities improve traceability, accountability,
operational efficiency, and stakeholder collaboration while supporting circular
economy objectives. Figure~\ref{fig:reverse-logistics} illustrates the reverse
logistics process adopted for electronic waste management, beginning with waste
generation and continuing through collection, transportation, inspection, repair,
refurbishment, recycling, and resource recovery.

\begin{figure}[H]
  \centering
  \imageplaceholder{Reverse Logistics Framework for Electronic Waste Management}
  \caption{Reverse Logistics Framework for Electronic Waste Management. Source: Adapted from Rogers and Tibben-Lembke (1999) and Govindan et al. (2015).}
  \label{fig:reverse-logistics}
\end{figure}

\subsection{Intelligent Waste Management}

Intelligent waste management refers to the application of modern information and
communication technologies to improve the planning, monitoring, collection,
transportation, processing, and disposal of waste materials. It integrates
technologies such as artificial intelligence, cloud computing, Internet of Things
(IoT), geographic information systems (GIS), mobile computing, data analytics, and
automation to enhance operational efficiency, decision-making, and environmental
sustainability \citep{belhadi2021,sharma2023}.

Traditional waste management systems largely depend on manual record
keeping, fixed collection schedules, and limited communication among
stakeholders. These approaches often result in inefficient resource utilization,
delayed waste collection, poor monitoring, and increased operational costs. In
contrast, intelligent waste management systems utilize digital technologies to
provide real-time monitoring, automated decision support, optimized route
planning, predictive analytics, and centralized information management.

The increasing availability of cloud computing and mobile technologies has
significantly accelerated the development of intelligent waste management
systems. Cloud platforms provide scalable storage and processing capabilities,
while mobile applications enable users to report waste, request collection services,
monitor collection progress, and receive real-time notifications. These
technologies improve communication between citizens, waste collectors, recycling
companies, and environmental authorities.

Artificial intelligence has further enhanced intelligent waste management
through automated waste classification, object detection, predictive analytics, and
decision support systems. Machine learning algorithms can analyse images of
waste materials, identify waste categories, estimate recyclable content, and
support automated sorting processes, thereby improving both efficiency and
accuracy \citep{abdallah2022}.

Geographic Information Systems (GIS) and Global Positioning System (GPS)
technologies also play important roles in intelligent waste management. These
technologies enable route optimization, location tracking, collection planning, and
spatial analysis, allowing organizations to reduce transportation costs while
improving service delivery and environmental performance.

Recent research indicates that integrating artificial intelligence, cloud
computing, GIS, and mobile technologies significantly improves waste collection
efficiency, operational transparency, resource utilization, and environmental
monitoring. These intelligent systems also support data-driven policy formulation
and sustainable urban development \citep{kaza2018,sharma2023}.

The EcoTrace system adopts the principles of intelligent waste management by
integrating cloud-based information management, artificial-intelligence-assisted
electronic waste classification, cross-platform mobile, web, and desktop
applications, Google Maps Platform for location services, Firebase Cloud
Messaging for real-time notifications, and centralized reporting dashboards.
Through these integrated technologies, the system improves electronic waste
traceability, operational efficiency, stakeholder collaboration, and decision-making
across the electronic waste management lifecycle. Figure~\ref{fig:intelligent-waste}
presents the major components of an intelligent waste management ecosystem and
illustrates how modern digital technologies support efficient waste management
operations.

\begin{figure}[H]
  \centering
  \imageplaceholder{Intelligent Waste Management Framework (AI, Cloud Computing, GIS, IoT, Mobile, Data Analytics)}
  \caption{Intelligent Waste Management Framework showing the integration of Artificial Intelligence, Cloud Computing, GIS, IoT, Mobile Applications, and Data Analytics in modern waste management. Source: Adapted from recent smart waste management literature \citep{sharma2023}.}
  \label{fig:intelligent-waste}
\end{figure}

\subsection{Cloud Computing}

Cloud computing is a computing paradigm that enables on-demand access to
configurable computing resources, including servers, storage, databases,
networking, software, and processing power, over the Internet. Instead of relying
on local infrastructure, cloud computing allows organizations to access scalable
computing services that can be provisioned rapidly with minimal management
effort \citep{armbrust2010,mell2011}.

Cloud computing has transformed the development and deployment of modern
information systems by providing scalable, flexible, and cost-effective computing
environments. Organizations increasingly adopt cloud technologies because they
reduce infrastructure costs, improve system availability, simplify maintenance, and
support collaboration through centralized data management.

The National Institute of Standards and Technology (NIST) identifies five
essential characteristics of cloud computing: on-demand self-service, broad
network access, resource pooling, rapid elasticity, and measured service
\citep{mell2011}. These characteristics enable cloud platforms to dynamically
allocate computing resources according to user demand while ensuring efficient
resource utilization.

Cloud computing services are commonly classified into three service models:

\begin{itemize}
  \item \textbf{Infrastructure as a Service (IaaS)}, which provides virtualized
    computing resources such as servers, storage, and networking.
  \item \textbf{Platform as a Service (PaaS)}, which provides software
    development platforms and deployment environments without requiring
    developers to manage the underlying infrastructure.
  \item \textbf{Software as a Service (SaaS)}, which delivers complete software
    applications over the Internet for end users.
\end{itemize}

Cloud computing also supports different deployment models, including public
cloud, private cloud, hybrid cloud, and community cloud. Public cloud services are
the most widely adopted because they offer scalability, accessibility, and lower
operational costs for organizations of different sizes.

Recent advances in cloud computing have enabled the integration of mobile
computing, artificial intelligence, Internet of Things (IoT), big data analytics, and
geographic information systems into modern software solutions. These
technologies facilitate real-time data processing, secure information sharing,
centralized storage, and intelligent decision support, making cloud computing an
essential component of smart environmental management systems
\citep{marinescu2022}.

The EcoTrace system extensively utilizes cloud computing to provide a secure,
scalable, and reliable electronic waste management platform. Cloud Firestore is
used as the primary cloud database for storing and synchronizing application data
in real time, Firebase Authentication provides secure user identity management,
Firebase Cloud Messaging delivers push notifications, Cloudinary manages
cloud-based image storage, and the FastAPI backend is deployed on the Render
cloud platform. Together, these cloud services enable centralized data
management, cross-platform accessibility, real-time communication, and high
system availability.

The adoption of cloud computing within EcoTrace enhances scalability, improves
data accessibility, supports collaboration among multiple stakeholders, and
simplifies system maintenance while reducing the need for organizations to
manage their own physical server infrastructure. Figure~\ref{fig:cloud-arch}
illustrates the role of cloud computing in supporting centralized data management
and communication between multiple client applications and cloud services.

\begin{figure}[H]
  \centering
  \imageplaceholder{Cloud Computing Architecture for Intelligent Electronic Waste Management}
  \caption{Cloud Computing Architecture for Intelligent Electronic Waste Management. Source: Adapted from Mell and Grance (2011) and recent cloud computing literature.}
  \label{fig:cloud-arch}
\end{figure}

\subsection{Artificial Intelligence}

Artificial Intelligence (AI) refers to the branch of computer science concerned with
developing computer systems capable of performing tasks that normally require
human intelligence. Such tasks include learning, reasoning, problem-solving,
pattern recognition, decision-making, speech recognition, natural language
processing, and computer vision \citep{russell2021}.

Artificial intelligence has become one of the most transformative technologies
of the Fourth Industrial Revolution. Advances in computing power, cloud
computing, big data, and machine learning algorithms have enabled AI systems to
solve complex real-world problems across numerous domains, including
healthcare, agriculture, transportation, manufacturing, education, finance, and
environmental management.

Machine Learning (ML), a major subfield of artificial intelligence, enables
computer systems to learn patterns directly from data without being explicitly
programmed for every possible scenario. Instead, machine learning algorithms
improve their performance by analysing historical data and identifying
relationships that support accurate predictions and classifications.

Deep Learning, a specialized branch of machine learning, utilizes Artificial
Neural Networks (ANNs) with multiple hidden layers to learn complex
representations from large datasets. Deep learning has achieved remarkable
success in image classification, object detection, speech recognition, autonomous
vehicles, and medical diagnosis due to its ability to automatically extract
meaningful features from raw data \citep{goodfellow2016}.

Computer vision, another important area of artificial intelligence, focuses on
enabling computers to interpret and analyse digital images and videos. Modern
computer vision systems employ deep learning models to identify objects, classify
images, detect anomalies, and perform image segmentation with high accuracy.
These capabilities have made computer vision an essential technology for
intelligent waste management, where waste materials can be automatically
identified and classified based on uploaded images.

Recent studies demonstrate that artificial intelligence significantly improves
waste management by automating waste classification, supporting smart
recycling, optimizing collection operations, predicting waste generation, and
enhancing environmental monitoring \citep{abdallah2022,sharma2023}. AI-assisted
systems reduce manual sorting errors, improve operational efficiency, and
increase the recovery of recyclable materials.

The effectiveness of artificial intelligence systems largely depends on the
quality and diversity of training datasets. Well-labelled and representative
datasets improve the ability of machine learning models to generalize effectively
when presented with previously unseen data. Conversely, poor-quality datasets
may reduce model accuracy and introduce classification bias.

Within the EcoTrace system, artificial intelligence is employed to support the
automatic classification of electronic waste images uploaded by users. The AI
component assists users in identifying the category of electronic waste before
collection, thereby improving data quality and supporting more efficient
collection, repair, refurbishment, and recycling processes. PyTorch was selected as
the machine learning framework because of its flexibility, extensive deep learning
support, active research community, and suitability for computer vision
applications.

The integration of artificial intelligence into EcoTrace demonstrates how
intelligent software systems can improve environmental management by
supporting faster decision-making, increasing classification accuracy, and
promoting sustainable electronic waste recovery practices. Figure~\ref{fig:ai-framework}
illustrates the general workflow of an artificial intelligence image classification
system used for electronic waste management.

\begin{figure}[H]
  \centering
  \imageplaceholder{Artificial Intelligence Framework for Electronic Waste Image Classification}
  \caption{Artificial Intelligence Framework for Electronic Waste Image Classification. Source: Adapted from Goodfellow et al. (2016) and Russell and Norvig (2021).}
  \label{fig:ai-framework}
\end{figure}

\subsection{Geographic Information Systems (GIS)}

Geographic Information Systems (GIS) are computer-based systems used for
capturing, storing, managing, analysing, and visualizing geographically referenced
data. GIS integrates spatial information with attribute data to support
decision-making through digital maps, spatial analysis, and location-based
services \citep{longley2021}.

The increasing adoption of GIS technologies has transformed numerous
industries, including urban planning, transportation, agriculture, healthcare,
disaster management, environmental monitoring, and waste management. By
combining geographic coordinates with real-time data, GIS enables organizations
to monitor assets, analyse spatial patterns, optimize operations, and improve
resource allocation.

Within waste management, GIS plays an important role in planning waste
collection routes, identifying waste generation hotspots, locating collection
centres, monitoring vehicle movement, and supporting environmental planning.
Spatial analysis enables waste management organizations to improve operational
efficiency while reducing fuel consumption, travel distance, collection time, and
operational costs \citep{ananth2022}.

Modern GIS applications are frequently integrated with Global Positioning
System (GPS) technologies to provide real-time location tracking. GPS enables
mobile devices and vehicles to determine their geographic positions accurately,
while GIS visualizes and analyses these locations to support intelligent
decision-making. The integration of GIS and GPS has become essential for smart
cities and intelligent waste management systems.

Cloud computing has further enhanced GIS capabilities by enabling cloud-based
map services, real-time spatial data synchronization, and cross-platform
accessibility. Cloud GIS solutions allow multiple users to access and update
geographic information simultaneously from mobile, web, and desktop
applications, thereby improving collaboration among stakeholders.

The Google Maps Platform is one of the most widely adopted cloud-based GIS
solutions for software development. It provides Application Programming
Interfaces (APIs) and software development kits (SDKs) that support interactive
maps, geocoding, reverse geocoding, route planning, distance calculation, place
search, and navigation services \citep{google2026maps}. These capabilities have
made Google Maps Platform a preferred solution for location-aware applications.

Several recent studies have demonstrated that integrating GIS into waste
management systems significantly improves collection efficiency, route
optimization, environmental monitoring, and service delivery
\citep{ananth2022,belhadi2021}. GIS also supports evidence-based decision-making
by enabling administrators to visualize operational activities and identify areas
requiring intervention.

Within the EcoTrace system, GIS functionality is implemented through the
Google Maps Platform to support electronic waste collection and transportation.
Users can register the geographic location of electronic waste, view nearby
collection centres, visualize collection routes, and track collection activities.
Environmental officers and administrators can monitor collection operations using
interactive maps, thereby improving operational coordination and
decision-making.

The integration of GIS within EcoTrace enhances route planning, reduces
transportation costs, improves service delivery, strengthens electronic waste
traceability, and contributes to more efficient and sustainable electronic waste
management. Figure~\ref{fig:gis-framework} illustrates how Geographic
Information Systems support location-based electronic waste management
through mapping, GPS tracking, route optimization, and spatial analytics.

\begin{figure}[H]
  \centering
  \imageplaceholder{GIS Framework for Intelligent Electronic Waste Management}
  \caption{Geographic Information System (GIS) Framework for Intelligent Electronic Waste Management. Source: Adapted from Longley et al. (2021) and Google Maps Platform documentation.}
  \label{fig:gis-framework}
\end{figure}

\section{Review of Relevant Technologies}

The successful development of modern intelligent information systems depends
on selecting technologies that provide reliability, scalability, security,
maintainability, and high performance. The development of EcoTrace required the
integration of several modern software technologies to support cross-platform
application development, cloud computing, secure authentication, image
management, real-time notifications, geographic information systems, and
artificial intelligence.

This section reviews the major technologies adopted in the implementation of
EcoTrace. For each technology, its features, advantages, limitations, and relevance
to the proposed system are discussed. The review provides the technological
justification for the selection of Flutter, FastAPI, Cloud Firestore, Firebase
Authentication, Firebase Cloud Messaging, Cloudinary, Google Maps Platform,
Render, and PyTorch.

\subsection{Flutter Framework}

Flutter is Google's open-source software development framework for building
cross-platform applications from a single codebase \citep{google2025flutter}. Since
its introduction in 2018, Flutter has become one of the most widely adopted
cross-platform development frameworks because it enables developers to build
applications for Android, iOS, Web, Windows, macOS, and Linux while
maintaining a consistent user interface and user experience.

Flutter is built on the Dart programming language and utilizes its own
high-performance rendering engine known as Skia. Unlike many cross-platform
frameworks that rely on native user interface widgets, Flutter renders its own
widgets, enabling consistent application appearance across different operating
systems.

One of Flutter's major advantages is the concept of a single codebase.
Developers write the application once and deploy it to multiple platforms with
minimal platform-specific modifications. This approach reduces development time,
maintenance costs, and software complexity while improving consistency across
supported devices.

Flutter also provides Hot Reload functionality, allowing developers to view code
changes almost instantly without restarting the application. This significantly
improves productivity during software development and debugging.

The framework includes an extensive collection of customizable widgets,
responsive layouts, animation libraries, navigation components, and integration
support for cloud services, RESTful APIs, databases, mapping services, and
artificial intelligence frameworks. These capabilities make Flutter suitable for
developing enterprise-scale applications.

Several comparative studies have reported that Flutter provides better
development productivity, faster application performance, and improved
maintainability compared to many traditional cross-platform frameworks
\citep{biswas2023}.

For the EcoTrace project, Flutter was selected because it supports the
development of Android mobile applications, Web applications, and Windows
desktop applications using a unified codebase. This approach ensures feature
consistency across all supported platforms while reducing development effort and
simplifying future system maintenance. Flutter also integrates seamlessly with
Firebase services, Google Maps Platform, Cloudinary, and RESTful APIs
developed using FastAPI, making it well suited to the requirements of the
proposed system. Figure~\ref{fig:flutter-arch} illustrates Flutter's cross-platform
architecture and its integration with the major cloud services used by EcoTrace.

\begin{figure}[H]
  \centering
  \imageplaceholder{Flutter Cross-Platform Architecture for EcoTrace}
  \caption{Flutter Cross-Platform Architecture for EcoTrace. Source: Adapted from the official Flutter documentation \citep{google2025flutter}.}
  \label{fig:flutter-arch}
\end{figure}

\subsection{FastAPI}

FastAPI is a modern, high-performance Python web framework used for building
Application Programming Interfaces (APIs). Developed by Sebasti\'an Ram\'irez in
2018, FastAPI is designed to simplify backend development while providing
exceptional performance, automatic API documentation, asynchronous request
handling, and built-in data validation \citep{ramirez2025}.

FastAPI is built on the ASGI (Asynchronous Server Gateway Interface) standard
and utilizes Starlette for web functionality and Pydantic for data validation. This
architecture enables FastAPI applications to handle multiple client requests
efficiently while maintaining excellent performance and scalability. Benchmark
studies have shown that FastAPI achieves performance levels comparable to
Node.js and Go for many web service workloads while offering the simplicity and
readability of the Python programming language \citep{ramirez2023}.

One of the major advantages of FastAPI is its automatic generation of
interactive API documentation using the OpenAPI Specification. As soon as an API
is implemented, FastAPI automatically generates Swagger UI and ReDoc
documentation, allowing developers to visualize endpoints, test API requests, and
simplify integration between frontend and backend systems.

FastAPI also provides robust support for asynchronous programming using
Python's async and await keywords. Asynchronous request processing allows the
server to handle numerous simultaneous client requests efficiently without
blocking execution, thereby improving system responsiveness and resource
utilization.

Another important feature of FastAPI is its integration with Pydantic models for
request validation and serialization. Incoming data is validated automatically
against predefined schemas before being processed by the application, reducing
programming errors and improving overall system reliability.

Security is also an important consideration in FastAPI applications. The
framework supports OAuth2 authentication, JSON Web Tokens (JWT), API key
authentication, Cross-Origin Resource Sharing (CORS), dependency injection,
request validation, and middleware integration, making it suitable for
enterprise-level web applications.

FastAPI integrates effectively with cloud services, databases, machine learning
frameworks, and third-party APIs. These capabilities have made it increasingly
popular for developing backend services supporting artificial intelligence,
cloud-native applications, Internet of Things (IoT), and microservice architectures.

Within the EcoTrace system, FastAPI serves as the primary backend framework
responsible for implementing RESTful API services that facilitate communication
between Flutter client applications and cloud services. The backend manages user
authentication requests, electronic waste registration, collection scheduling, image
processing, notification services, reporting, and communication with Cloud
Firestore, Cloudinary, Google Maps Platform, Firebase Cloud Messaging, and the AI
image classification module.

The FastAPI backend is deployed on the Render cloud platform, enabling
secure, scalable, and reliable API hosting. This deployment supports cross-platform
communication between Android, Web, and Windows applications while providing
centralized business logic and secure access to cloud resources.

Overall, FastAPI provides the performance, scalability, maintainability, and
developer productivity required for implementing a modern cloud-based
electronic waste management platform such as EcoTrace. Figure~\ref{fig:fastapi-arch}
illustrates the architecture of the FastAPI backend and its interaction with the
various cloud services and client applications used by EcoTrace.

\begin{figure}[H]
  \centering
  \imageplaceholder{FastAPI Backend Architecture for EcoTrace}
  \caption{FastAPI Backend Architecture for EcoTrace. Source: Adapted from the official FastAPI documentation \citep{ramirez2025}.}
  \label{fig:fastapi-arch}
\end{figure}

\subsection{Cloud Firestore}

Cloud Firestore is a fully managed, cloud-hosted NoSQL document database
provided by Google as part of the Firebase platform. It is designed to store,
synchronize, and retrieve application data in real-time while providing high
scalability, reliability, and security \citep{firebase2025firestore}.

Unlike traditional relational databases that organize information into tables
consisting of rows and columns, Cloud Firestore stores data as collections of
documents. Each document contains fields represented as key-value pairs and can
itself contain nested collections. This document-oriented architecture provides
greater flexibility for representing complex and hierarchical application data.

One of the major advantages of Cloud Firestore is its real-time data
synchronization capability. Whenever data is added, modified, or deleted, changes
are automatically synchronized across all connected client applications without
requiring manual refresh operations. This feature is particularly useful for
applications that require multiple users to access current information
simultaneously.

Cloud Firestore also supports offline data persistence. Client applications can
continue reading and writing data while temporarily disconnected from the
Internet. Once connectivity is restored, local changes are automatically
synchronized with the cloud database, improving reliability and user experience in
environments with unstable network connections.

Another important characteristic of Cloud Firestore is its automatic horizontal
scalability. Google manages infrastructure provisioning, load balancing,
replication, backup, and disaster recovery, allowing developers to focus on
application development instead of database administration. This makes Cloud
Firestore suitable for applications expected to grow in both user base and data
volume.

Security within Cloud Firestore is enforced through Firebase Security Rules,
which allow developers to define fine-grained access policies based on
authenticated users, user roles, document ownership, and other conditions. These
rules ensure that only authorized users can access or modify sensitive application
data.

Cloud Firestore integrates seamlessly with other Firebase services, including
Firebase Authentication, Firebase Cloud Messaging, Cloud Functions, Firebase
Storage, Analytics, and Google Cloud services. This close integration simplifies the
development of cloud-native applications while improving maintainability and
interoperability.

Several recent studies have identified cloud-hosted NoSQL databases as
effective solutions for scalable mobile and web applications because they support
flexible data structures, real-time synchronization, and distributed cloud
architectures \citep{li2022,nayak2021}.

Within the EcoTrace system, Cloud Firestore serves as the primary cloud
database responsible for storing user profiles, electronic waste records, pickup
requests, collection schedules, recycling information, notifications, reports, and
administrative data. Because Android, Web, and Windows applications access the
same centralized database, authorized users can retrieve consistent information
regardless of the platform being used.

The adoption of Cloud Firestore contributes significantly to the scalability,
reliability, availability, and real-time communication capabilities of EcoTrace while
simplifying database administration and supporting future system expansion.
Figure~\ref{fig:firestore-arch} illustrates the Cloud Firestore architecture adopted
by EcoTrace for centralized data storage and real-time synchronization.

\begin{figure}[H]
  \centering
  \imageplaceholder{Cloud Firestore Architecture for EcoTrace}
  \caption{Cloud Firestore Architecture for EcoTrace. Source: Adapted from Firebase Documentation \citep{firebase2025firestore}.}
  \label{fig:firestore-arch}
\end{figure}

\subsection{Firebase Authentication}

Firebase Authentication is a cloud-based identity management service provided by
Google as part of the Firebase platform. It simplifies user authentication by
providing secure mechanisms for user registration, login, password management,
email verification, and session management without requiring developers to
implement complex authentication infrastructure \citep{firebase2025auth}.

Authentication is a fundamental security requirement for modern information
systems because it ensures that only authorized users can access protected
resources and perform system operations. Firebase Authentication supports
multiple authentication methods, including email and password authentication,
phone number authentication, Google Sign-In, Apple Sign-In, Facebook Login,
GitHub authentication, Microsoft authentication, anonymous authentication, and
custom authentication providers.

One of the major advantages of Firebase Authentication is its seamless
integration with other Firebase services, including Cloud Firestore, Firebase Cloud
Messaging, Cloud Functions, Firebase Storage, and Google Cloud services. This
integration simplifies user identity management while maintaining secure
communication across cloud-based applications.

Firebase Authentication implements modern security standards, including
encrypted communication using HTTPS, secure password storage, token-based
authentication, session management, and JSON Web Tokens (JWT). These
mechanisms protect user identities and reduce the risk of unauthorized system
access.

Role-Based Access Control (RBAC) is commonly implemented alongside
Firebase Authentication to restrict user access according to assigned roles and
responsibilities. Under the RBAC model, users are granted permissions based on
predefined roles rather than individual user accounts. This approach simplifies
security administration and ensures that users can only access functionalities
relevant to their assigned responsibilities \citep{sandhu1996}.

The EcoTrace system adopts Firebase Authentication together with Role-Based
Access Control to manage multiple categories of users, including households,
businesses, waste collectors, technicians, recyclers, environmental officers,
administrators, and super administrators. Each authenticated user is granted
specific permissions according to their assigned role, thereby protecting sensitive
information and preventing unauthorized access to system resources.

Firebase Authentication also supports persistent user sessions, allowing users
to remain logged in across application restarts while maintaining secure
authentication tokens. Password reset, email verification, and account recovery
mechanisms further improve system security and usability.

The adoption of Firebase Authentication significantly improves the security,
reliability, and maintainability of EcoTrace by providing a trusted cloud-based
authentication infrastructure while reducing the development effort required to
implement secure user identity management. Figure~\ref{fig:firebase-auth-arch}
illustrates the authentication workflow implemented within the EcoTrace system.

\begin{figure}[H]
  \centering
  \imageplaceholder{Firebase Authentication Architecture for EcoTrace}
  \caption{Firebase Authentication Architecture for EcoTrace. Source: Adapted from the official Firebase Authentication documentation \citep{firebase2025auth}.}
  \label{fig:firebase-auth-arch}
\end{figure}

\subsection{Firebase Cloud Messaging (FCM)}

Firebase Cloud Messaging (FCM) is Google's cloud-based messaging service that
enables developers to send real-time notifications and data messages to Android,
Web, and other supported client applications. It provides a reliable
communication channel between backend services and client devices without
requiring developers to manage their own notification servers \citep{firebase2025fcm}.

Modern information systems increasingly depend on real-time communication
to improve user engagement, operational efficiency, and service delivery. Push
notification services enable organizations to instantly communicate important
events, reminders, alerts, and system updates to users regardless of whether the
application is currently active.

Firebase Cloud Messaging supports two primary message types: notification
messages and data messages. Notification messages are displayed directly by the
operating system, while data messages are processed by the application itself,
allowing developers to implement custom notification behaviour and business
logic.

One of the major advantages of Firebase Cloud Messaging is its seamless
integration with Firebase Authentication, Cloud Firestore, Cloud Functions, and
mobile application frameworks such as Flutter. This integration simplifies the
implementation of cloud-based notification systems while ensuring secure and
reliable message delivery.

FCM provides several features that improve communication efficiency,
including topic messaging, device group messaging, targeted user notifications,
scheduled notifications, and message prioritization. These capabilities enable
organizations to deliver personalized messages to individual users, specific user
groups, or all users of an application.

Cloud-based notification services have become increasingly important for
intelligent environmental management systems because they improve coordination
among stakeholders and facilitate timely decision-making. Real-time alerts help
reduce operational delays, improve user participation, and enhance service quality.

Within the EcoTrace system, Firebase Cloud Messaging is used to send
real-time notifications throughout the electronic waste management process.
Users receive notifications when pickup requests are submitted, approved,
scheduled, completed, or cancelled. Additional notifications inform users about
recycling progress, collection reminders, environmental campaigns, system
announcements, and account activities.

Environmental officers, collectors, recyclers, technicians, and administrators
also receive operational notifications regarding new assignments, collection
requests, schedule updates, and system activities. These notifications improve
coordination among system stakeholders while reducing communication delays.

The integration of Firebase Cloud Messaging significantly enhances the overall
usability of EcoTrace by improving user engagement, strengthening
communication, supporting timely service delivery, and providing a responsive
user experience across Android, Web, and Windows platforms.
Figure~\ref{fig:fcm-arch} illustrates the notification workflow implemented in
EcoTrace using Firebase Cloud Messaging.

\begin{figure}[H]
  \centering
  \imageplaceholder{Firebase Cloud Messaging Architecture for EcoTrace}
  \caption{Firebase Cloud Messaging Architecture for EcoTrace. Source: Adapted from the official Firebase Cloud Messaging documentation \citep{firebase2025fcm}.}
  \label{fig:fcm-arch}
\end{figure}

\subsection{Cloudinary}

Cloudinary is a cloud-based media management platform that provides services for
storing, managing, optimizing, transforming, and delivering images and videos
through cloud infrastructure. It enables developers to upload media files securely
while automatically optimizing them for different devices, browsers, and network
conditions \citep{cloudinary2025}.

Modern software applications increasingly depend on cloud-based media
management services because digital images constitute a significant portion of
application data. Traditional file storage approaches often require developers to
manually implement image compression, resizing, format conversion, caching,
and content delivery mechanisms. Cloudinary automates these processes, thereby
improving application performance while reducing development complexity.

One of Cloudinary's major strengths is its automatic image optimization
capability. Uploaded images are automatically compressed, resized, cropped, and
converted into modern formats such as WebP or AVIF where supported, reducing
bandwidth consumption while maintaining high visual quality. These
optimizations contribute to faster page loading times and improved user
experience.

Cloudinary also integrates a global Content Delivery Network (CDN), allowing
media assets to be served from geographically distributed servers. This reduces
image loading latency by delivering media from servers located closer to end
users. Such capabilities are especially important for cloud-based applications that
serve users across different geographic regions.

Another important feature of Cloudinary is its support for secure media
management. Developers can configure authenticated uploads, signed URLs,
access control policies, and transformation restrictions to protect media assets
from unauthorized access. Cloudinary also provides automatic backups, version
control, and high availability through its managed cloud infrastructure.

Cloudinary offers comprehensive Application Programming Interfaces (APIs)
and Software Development Kits (SDKs) for numerous programming languages and
frameworks, including Python, JavaScript, Flutter, and Node.js. These APIs
simplify media upload, retrieval, transformation, and deletion while supporting
integration with cloud-native applications.

Within the EcoTrace system, Cloudinary is responsible for storing electronic
waste images uploaded by users during waste registration. Each uploaded image
is securely stored in the cloud, optimized for delivery, and linked to the
corresponding waste record stored in Cloud Firestore. This separation of
structured data and media storage improves scalability and reduces database
storage requirements.

The backend manages communication with the Cloudinary API, performing
secure image uploads and storing the returned image URLs in Cloud Firestore.
Flutter client applications subsequently retrieve these optimized image URLs
when displaying electronic waste records, thereby reducing application size and
improving loading performance.

The adoption of Cloudinary significantly enhances EcoTrace by providing
reliable cloud-based image hosting, automatic media optimization, secure storage,
high availability, and efficient image delivery across Android, Web, and Windows
applications. Figure~\ref{fig:cloudinary-arch} illustrates how Cloudinary integrates
with the EcoTrace architecture to manage electronic waste images.

\begin{figure}[H]
  \centering
  \imageplaceholder{Cloudinary Architecture for Image Management in EcoTrace}
  \caption{Cloudinary Architecture for Image Management in EcoTrace. Source: Adapted from the official Cloudinary documentation \citep{cloudinary2025}.}
  \label{fig:cloudinary-arch}
\end{figure}

\subsection{Google Maps Platform}

Google Maps Platform is a collection of cloud-based mapping, geospatial, routing,
navigation, and location services provided by Google. It offers Application
Programming Interfaces (APIs) and Software Development Kits (SDKs) that enable
developers to integrate interactive maps and location-aware functionality into
mobile, web, desktop, and server-side applications \citep{google2026maps}.

Location-based services are important components of modern logistics,
transportation, environmental monitoring, and waste management applications.
They enable organizations to record geographic coordinates, display operational
locations, calculate travel distances, plan routes, locate nearby facilities, and
monitor the spatial distribution of activities.

The Google Maps Platform includes several services that are relevant to the
EcoTrace system. The Maps services support the display of interactive maps,
markers, routes, and geographic features. The Geocoding API converts
human-readable addresses into latitude and longitude coordinates, while reverse
geocoding converts geographic coordinates into readable addresses
\citep{google2026geocoding}.

The Places API provides information about geographic locations and facilities.
It supports place search, place details, autocomplete, nearby search, and the use
of unique place identifiers. These capabilities can assist users in finding electronic
waste collection centres and identifying relevant facilities near their current
locations \citep{google2026places}.

The Routes API supports route planning between origins, destinations, and
intermediate waypoints. Its Compute Routes operation determines an appropriate
route between locations, while Compute Route Matrix calculates travel distances
and estimated journey times for multiple origins and destinations
\citep{google2026routes}.

Google Maps Platform also provides route optimization capabilities. Waypoint
optimization can rearrange intermediate stops to improve travel efficiency, while
the Route Optimization API can assign collection tasks and construct efficient
routes for vehicle fleets according to objectives and operational constraints
\citep{google2026optway,google2026routeopt}.

These routing services are particularly relevant to electronic waste
management because collection requests may originate from multiple geographic
locations. Efficient route planning can reduce unnecessary travel, fuel
consumption, collection delays, and transportation costs while improving service
coverage.

Within EcoTrace, Google Maps Platform is used to display electronic waste
locations, collection centres, pickup routes, and transportation activities.
Geographic coordinates associated with pickup requests are processed through
the backend and stored with the relevant records in Cloud Firestore.

Collectors and drivers can use the generated route information to navigate
between waste collection locations and processing facilities. Administrators and
environmental officers can also use map-based dashboards to monitor collection
coverage, operational locations, and transportation activities.

Google Maps Platform requires authenticated API access and careful credential
management. API keys should be restricted to the specific applications and
services that require them to reduce unauthorized use and unexpected billing
\citep{google2026bestpractices}. The platform also operates primarily under
usage-based pricing, which means that API consumption should be monitored
and controlled through quotas and budget alerts \citep{google2026pricing}.

The adoption of Google Maps Platform improves EcoTrace by providing
accurate location visualization, route planning, proximity analysis,
collection-centre discovery, and route optimization. These capabilities strengthen
operational coordination and support more efficient and sustainable electronic
waste collection. Figure~\ref{fig:maps-arch} illustrates how Google Maps Platform
supports location registration, collection-centre discovery, route calculation, and
route optimization within EcoTrace.

\begin{figure}[H]
  \centering
  \imageplaceholder{Google Maps Platform Architecture for EcoTrace}
  \caption{Google Maps Platform Architecture for EcoTrace. Source: Adapted from Google Maps Platform documentation \citep{google2026maps,google2026routes}.}
  \label{fig:maps-arch}
\end{figure}

\subsection{Render Cloud Platform}

Render is a cloud application hosting platform that enables developers to deploy,
operate, and monitor web applications, backend services, Application
Programming Interfaces (APIs), static websites, databases, and other cloud
resources. The platform supports several programming languages and
frameworks, including Python and FastAPI, and provides managed infrastructure
that reduces the need for developers to configure and maintain physical servers
\citep{render2026webservices}.

Modern software applications require reliable hosting environments that
support deployment automation, secure communication, configuration
management, monitoring, and scalability. Render addresses these requirements by
allowing developers to connect application repositories hosted on GitHub, GitLab,
or Bitbucket and deploy them as cloud services.

A Render Web Service is suitable for applications that execute server-side code
and respond to HTTP requests. Each deployed service receives a publicly
accessible \texttt{onrender.com} address, while custom domains can also be
configured. Render automatically manages HTTPS connections and redirects
incoming HTTP traffic to HTTPS, supporting secure communication between client
applications and backend services \citep{render2026webservices}.

Render supports automatic deployment from a linked Git branch. Whenever
new code is pushed to the configured branch, the platform can build and deploy
the updated application automatically. This capability supports continuous
deployment and reduces the manual effort required to publish software updates.

For Python applications, Render executes a build command to install project
dependencies and a start command to launch the application. A FastAPI service
can be deployed using a build command such as \texttt{pip install -r
requirements.txt} and started with \texttt{uvicorn main:app --host 0.0.0.0
--port \$PORT}. The application must listen on the host \texttt{0.0.0.0} and use
the port supplied through Render's \texttt{PORT} environment variable
\citep{render2026fastapi}.

Render provides environment-variable and secret-file management for
protecting sensitive configuration information. API credentials, Firebase
service-account settings, Cloudinary credentials, application configuration values,
and other secrets can therefore be stored outside the source code and supplied
securely to the running service \citep{render2026envvars}.

The platform also supports health checks for monitoring application
availability. Render can periodically test a configured endpoint to confirm that a
service is ready to receive traffic. Unresponsive instances can be identified and
restarted, while newly deployed versions are verified before receiving production
traffic \citep{render2026healthchecks}.

Application and deployment logs are available through the Render Dashboard.
These logs assist developers in identifying startup errors, dependency problems,
failed requests, configuration issues, and other runtime events during
development and maintenance.

Within EcoTrace, Render hosts the backend that provides the central business
logic and RESTful API services. Flutter applications running on Android, Web, and
Windows communicate with the deployed API through secure HTTPS requests.

The Render-hosted backend processes electronic waste registration, image
upload, AI classification, pickup scheduling, notification requests, reporting, and
administrative operations. It also coordinates communication with Cloud
Firestore, Firebase Authentication, Firebase Cloud Messaging, Cloudinary, Google
Maps Platform, and the PyTorch AI component.

The EcoTrace backend currently uses the Render Free Web Service tier. This
tier is suitable for development, academic demonstration, and limited testing;
however, it introduces operational constraints. A free service is suspended after
fifteen minutes without inbound traffic and must restart when a new request is
received. The restart process may produce an initial response delay
\citep{render2026free}.

Free Web Services also use an ephemeral local filesystem. Files created or
uploaded directly to the server may be lost when the service restarts, redeploys, or
spins down. EcoTrace therefore stores structured application data in Cloud
Firestore and hosts uploaded images in Cloudinary rather than relying on the
Render filesystem.

The free tier also provides limited compute resources and does not support
horizontal scaling beyond a single service instance. These limitations may affect
response time when processing computationally intensive AI inference requests or
handling a large number of simultaneous users. A paid or dedicated hosting
configuration would therefore be more appropriate for large-scale production
deployment.

Despite these limitations, Render was selected because it provides a
straightforward deployment process, Git-based automation, managed HTTPS,
environment-variable management, logging, health checks, and direct support for
Python and FastAPI. These capabilities make it suitable for hosting the EcoTrace
backend during development, testing, demonstration, and initial deployment.
Figure~\ref{fig:render-arch} illustrates the deployment of the EcoTrace backend on
the Render cloud platform and its connections to the client applications and
external cloud services.

\begin{figure}[H]
  \centering
  \imageplaceholder{Render Deployment Architecture for EcoTrace}
  \caption{Render Deployment Architecture for EcoTrace. Source: Adapted from the official Render documentation \citep{render2026fastapi,render2026webservices}.}
  \label{fig:render-arch}
\end{figure}

\subsection{PyTorch Framework}

PyTorch is an open-source machine learning framework used for building,
training, evaluating, and deploying artificial intelligence models. It provides
optimized tensor operations for Central Processing Units (CPUs) and Graphics
Processing Units (GPUs), making it suitable for computationally intensive deep
learning applications \citep{paszke2019,pytorch2026docs}.

PyTorch follows an imperative and Python-oriented programming approach.
Operations are executed immediately, allowing developers to inspect intermediate
results, identify errors, and modify model behaviour during development. This
approach improves experimentation and debugging while maintaining support for
efficient computation and hardware acceleration \citep{paszke2019}.

A major component of PyTorch is its automatic differentiation system, known
as autograd. During model training, autograd records tensor operations and
automatically calculates gradients during backpropagation. These gradients are
used by optimization algorithms to update model parameters and progressively
reduce prediction error \citep{pytorch2026autograd}.

PyTorch provides several modules that support the machine learning
development lifecycle. The \texttt{torch.nn} module contains layers, activation
functions, loss functions, and other components used to construct neural
networks. The \texttt{torch.optim} module provides optimization algorithms,
while \texttt{Dataset} and \texttt{DataLoader} utilities support structured loading,
batching, and shuffling of training data.

For computer vision applications, PyTorch provides the TorchVision library.
TorchVision contains common image datasets, model architectures, pretrained
weights, and image transformation utilities. These resources support image
resizing, cropping, normalization, augmentation, training, validation, and
inference \citep{pytorch2026torchvision}.

Convolutional Neural Networks (CNNs) are widely used for image classification
because they learn spatial patterns such as edges, textures, shapes, and object
structures directly from images. Within an electronic waste classification system,
a CNN can analyse an uploaded image and predict whether it represents a mobile
phone, laptop, television, printer, battery, circuit board, or another supported
electronic waste category.

The performance of an image classification model depends on the quality,
quantity, diversity, and correctness of the training data. Images should represent
different device conditions, orientations, backgrounds, lighting conditions, camera
qualities, and levels of physical damage. Data augmentation can improve model
generalization by introducing controlled transformations such as rotation,
cropping, resizing, flipping, and brightness adjustment.

Transfer learning is particularly useful when the available project dataset is
not large enough to train a deep neural network completely from random
initialization. Under this approach, a model previously trained on a large image
dataset is adapted to the electronic waste classification task. The earlier layers
retain general visual features, while selected layers are fine-tuned to recognize
the required electronic waste categories \citep{chilamkurthy2025}.

A typical PyTorch image classification workflow consists of the following
stages: (1)~collection and labelling of electronic waste images; (2)~separation of
the dataset into training, validation, and testing subsets; (3)~image preprocessing
and augmentation; (4)~loading the images using datasets and data loaders;
(5)~selection and configuration of a suitable neural network; (6)~model training
and optimization; (7)~validation and hyperparameter adjustment; (8)~final testing
using previously unseen images; (9)~saving the trained model parameters; and
(10)~deployment of the trained model for inference.

Model performance may be evaluated using accuracy, precision, recall,
F1-score, a confusion matrix, and inference time. Accuracy measures the overall
proportion of correct predictions, while precision and recall provide more
detailed information about performance for individual electronic waste
categories. The confusion matrix helps identify categories that are frequently
misclassified.

PyTorch supports saving and restoring trained model parameters. The
recommended approach is to store the model's \texttt{state\_dict}, which contains
the learned parameters and can later be loaded into the same model architecture
for evaluation or inference \citep{pytorch2026saving}.

Within EcoTrace, PyTorch provides the artificial intelligence component
responsible for analysing electronic waste images. Images uploaded from the
Flutter applications are transmitted to the backend, which validates and
preprocesses them before passing them to the trained PyTorch model.

The model returns a predicted electronic waste category and a confidence
score. The backend formats the prediction as a JSON response, stores relevant
classification information in Cloud Firestore, and returns the result to the
Android, Web, or Windows application. Firebase Cloud Messaging may
subsequently notify the user when classification processing has been completed.

The exact neural network architecture for EcoTrace had not been finalized at
the time this chapter was written. Candidate architectures were to be evaluated
according to classification accuracy, precision, recall, F1-score, model size,
memory consumption, and inference time. This evaluation was considered
necessary because the selected model would need to provide acceptable
classification performance while operating within the computing constraints of
the deployed backend service. \textit{Chapter~\ref{ch:implementation} reports the
classification approach that was actually delivered.}

PyTorch was selected for EcoTrace because it provides a flexible Python-based
development environment, automatic differentiation, computer vision support,
pretrained model resources, hardware acceleration, and effective integration with
FastAPI. These capabilities make it suitable for developing and deploying the
electronic waste image classification component of the system.
Figure~\ref{fig:pytorch-arch} illustrates the training and inference workflows of
the PyTorch electronic waste classification component integrated into EcoTrace.

\begin{figure}[H]
  \centering
  \imageplaceholder{PyTorch Training and Inference Architecture for Electronic Waste Classification}
  \caption{PyTorch Training and Inference Architecture for Electronic Waste Classification in EcoTrace. Source: Adapted from PyTorch and TorchVision documentation \citep{pytorch2026docs,pytorch2026torchvision}.}
  \label{fig:pytorch-arch}
\end{figure}

\section{Review of Related and Existing Systems}

This section reviews selected electronic waste management systems and digital
platforms that demonstrate different approaches to electronic waste collection,
recycling, regulatory compliance, resource recovery, and circular economy
management. The systems were selected because their objectives or features are
related to one or more components of EcoTrace.

The review examines the Central Pollution Control Board E-Waste EPR Portal,
Recykal, ecoATM, E-Tadweer, and Green Grid. Each system is evaluated according
to its purpose, target users, principal functionalities, technological approach,
strengths, and limitations. The findings provide a foundation for the comparative
analysis presented in the next section.

\subsection{Central Pollution Control Board E-Waste EPR Portal}

The Central Pollution Control Board (CPCB) E-Waste Extended Producer
Responsibility (EPR) Portal is an online regulatory information system developed
to support the implementation of India's E-Waste Management Rules. The portal
provides a centralized environment in which producers, manufacturers, recyclers,
and refurbishers can register and manage their regulatory obligations
\citep{cpcb2026}.

The system supports producer registration, recycler registration, refurbisher
registration, submission of sales data, allocation of recycling targets, generation
and transfer of EPR certificates, annual returns, profile amendments, and
compliance monitoring. Through these features, the portal improves
documentation and accountability among stakeholders participating in the formal
electronic waste management sector.

A major strength of the CPCB portal is its regulatory orientation. It provides a
structured digital mechanism for monitoring compliance with national electronic
waste legislation and maintaining records of authorized stakeholders. It also
supports the implementation of EPR by linking producers' obligations with
recycling activities and certificates.

However, the publicly documented functions of the portal are primarily
focused on regulatory registration, certificates, targets, and returns. It is not
principally designed as a citizen-facing platform for electronic waste image
classification, household pickup scheduling, route optimization, repair
management, marketplace activities, or real-time lifecycle tracking.

EcoTrace differs from the CPCB portal by combining operational and
regulatory functions within a user-oriented cross-platform system. In addition to
administrative monitoring, EcoTrace supports households, businesses, collectors,
drivers, technicians, recyclers, and environmental officers throughout the
electronic waste recovery lifecycle.

\subsection{Recykal}

Recykal is a digital circular economy platform developed to connect waste
generators, aggregators, processors, recyclers, brands, and other stakeholders
involved in material recovery. The platform provides digital solutions for waste
transactions, Extended Producer Responsibility compliance, collection
management, traceability, and recycling operations \citep{recykal2026b}.

Its electronic waste disposal solution enables organizations to manage
obsolete information technology assets and arrange secure, sustainable, and
documented disposal. The platform emphasizes authorized processing,
traceability, compliance, data protection, and resource recovery
\citep{recykal2026a}.

A major strength of Recykal is its broad ecosystem approach. Rather than
supporting only one category of stakeholder, the platform connects multiple
participants across the circular economy value chain. Its digital records and
transaction mechanisms improve transparency and support organizations seeking
evidence of responsible waste processing.

Recykal is primarily a commercial circularity and compliance platform
designed for established waste and recycling markets. Its publicly documented
e-waste services place substantial emphasis on enterprise asset disposal,
compliance, traceability, and certified recycling. Consequently, it does not directly
reflect every feature required for the Sierra Leonean context, such as a unified
household application, collector and driver dashboards, repair workflow
management, integrated environmental-officer interfaces, and locally adaptable
operational roles.

EcoTrace adopts a similarly integrated circular-economy perspective but
extends the operational scope to include AI-assisted classification, pickup
scheduling, reverse logistics, repair, refurbishment, recycling, notifications,
location services, marketplace management, and environmental reporting across
Android, Web, and Windows applications.

\subsection{ecoATM}

ecoATM is an automated electronic device collection system that enables
consumers to sell eligible used devices through self-service kiosks. The kiosks are
placed in retail and other publicly accessible locations, providing a convenient
channel through which users can return unwanted mobile devices
\citep{ecoatm2026a}.

During a transaction, the user presents an eligible device and the required
identification. The kiosk examines the device, assesses its condition, generates a
price offer, and provides payment when the offer is accepted. Devices collected
through the system may subsequently be reused, resold, refurbished, or recycled
\citep{ecoatm2026b}.

The principal strength of ecoATM is convenience. Its automated kiosks reduce
the effort required for consumers to return unwanted devices and provide an
immediate financial incentive. The system also demonstrates the practical use of
automated inspection and valuation technologies in electronic device recovery.

Nevertheless, ecoATM concentrates primarily on selected consumer devices,
especially mobile phones and tablets, and depends on the availability of physical
kiosks. Its purpose is device acquisition and resale or recycling rather than
comprehensive management of multiple electronic waste categories.

The system does not provide the same broad operational workflow proposed
for EcoTrace, which includes household and business registration, pickup
scheduling, collector assignment, route management, repair, refurbishment,
recycling records, environmental monitoring, reporting, and role-specific
dashboards.

EcoTrace therefore extends beyond automated take-back by coordinating the
complete electronic waste recovery process through mobile, web, and Windows
applications while remaining adaptable to locations where specialized collection
kiosks are unavailable.

\subsection{E-Tadweer}

E-Tadweer is an electronic waste recycling application developed in Egypt to
encourage citizens to dispose of unwanted electrical and electronic equipment
through approved channels. The platform combines digital waste reporting,
collection-point information, and consumer incentives to improve participation in
formal electronic waste recycling \citep{etadweer2026}.

Users can submit information and images relating to electronic waste, identify
participating drop-off locations, and receive electronic vouchers or other rewards
from partner organizations. This incentive model encourages citizens to redirect
obsolete electronic devices away from informal or unsafe disposal channels.

A major strength of E-Tadweer is its focus on citizen participation. The
platform uses rewards to make responsible electronic waste disposal more
attractive while connecting users with formal collection arrangements. It also
demonstrates how mobile technology can improve public awareness and
participation in recycling programmes.

However, its publicly documented workflow concentrates mainly on consumer
submission, drop-off locations, rewards, and recycling participation. The available
information does not indicate the same level of integration for collector route
optimization, driver transportation records, repair and refurbishment
management, AI-assisted waste classification, environmental-officer monitoring,
or multi-platform administrative operations.

EcoTrace incorporates citizen participation but expands the lifecycle coverage
by supporting electronic waste registration, pickup requests, collection
assignments, transportation, repair, refurbishment, recycling, marketplace
operations, notifications, analytics, and environmental monitoring within one
integrated platform.

\subsection{Green Grid}

Green Grid is a recently proposed web-based electronic waste recycling platform
intended to improve public participation, awareness, collection access, and
recycling transparency. The system was developed as a full-stack web application
integrating several features associated with responsible electronic waste disposal
\citep{jagtap2026}.

Its reported functions include an electronic waste collection-point locator,
scheduled pickup services, a rewards mechanism, an awareness and educational
hub, a recycling-impact calculator, an AI assistant, and an eco-marketplace. The
platform uses web technologies, a relational database, mapping services, and
token-based authentication.

A major strength of Green Grid is its integrated user-engagement model. It
combines convenience, education, incentives, pickup scheduling, and
marketplace functionality rather than treating recycling as an isolated
transaction. The inclusion of mapping and AI-assisted features also demonstrates
the growing role of intelligent technologies in electronic waste management.

However, Green Grid is described primarily as a web-based platform and a
recent academic prototype. The available publication does not demonstrate a
fully deployed multi-platform operational system covering specialized
stakeholders such as drivers, technicians, recyclers, environmental officers,
administrators, and super administrators.

Its documented architecture also does not provide the same explicit
integration of mobile, web, and Windows applications with Cloud Firestore,
Firebase Authentication, Firebase Cloud Messaging, Cloudinary, Render-hosted
FastAPI services, and PyTorch-based image classification.

EcoTrace shares several objectives with Green Grid, including pickup services,
mapping, user engagement, AI support, and marketplace functionality. It extends
these concepts through role-specific applications, reverse-logistics tracking,
repair and refurbishment management, recycling records, notification services,
environmental analytics, and cross-platform deployment.
Figure~\ref{fig:related-systems} summarizes the principal orientation of the
electronic waste systems reviewed in this section and highlights the different
problems addressed by each platform.

\begin{figure}[H]
  \centering
  \imageplaceholder{Overview of Related Electronic Waste Management Systems}
  \caption{Overview of Related Electronic Waste Management Systems. Source: Author's compilation based on the reviewed systems.}
  \label{fig:related-systems}
\end{figure}

\section{Comparative Analysis of Existing Systems}

A comparative analysis was conducted to evaluate the functions, strengths, and
limitations of the electronic waste management systems reviewed in the
preceding section.

The analysis is based on publicly documented features. Therefore, the
expression \emph{Not documented} indicates that the reviewed sources did not
provide sufficient evidence that the relevant capability is available.

Table~\ref{tab:comparative-analysis} presents the comparative analysis of the
reviewed systems.

\begin{table}[H]
\centering
\caption{Comparative Analysis of Existing Electronic Waste Management Systems}
\label{tab:comparative-analysis}
\begin{tabular}{@{}p{2.6cm}p{2.6cm}p{1.6cm}p{1.6cm}p{2.2cm}p{2.6cm}@{}}
\toprule
\textbf{System} & \textbf{Primary Orientation} & \textbf{Pickup Sched.} & \textbf{AI Support} & \textbf{Location \& Routing} & \textbf{Repair, Marketplace, Compliance} \\
\midrule
CPCB EPR Portal & Regulatory compliance & Not doc. & Not doc. & Not doc. & Compliance only \\
Recykal & Enterprise circularity & Not doc. & Not doc. & Not doc. & Compliance, trading \\
ecoATM & Automated device take-back & Kiosk-based & Not doc. & Kiosk locations & Not doc. \\
E-Tadweer & Citizen participation \& rewards & Drop-off & Not doc. & Drop-off points & Rewards \\
Green Grid & Integrated web platform & Yes & AI assistant & Mapping & Marketplace \\
\textbf{EcoTrace} & \textbf{Integrated lifecycle platform} & \textbf{Yes} & \textbf{Yes (planned)} & \textbf{Yes (planned)} & \textbf{Repair, marketplace, reporting} \\
\bottomrule
\end{tabular}
\end{table}

\subsection{Discussion of the Comparison}

The comparison demonstrates that the reviewed systems address important
components of electronic waste management, but they differ considerably in their
objectives and functional coverage. The CPCB E-Waste EPR Portal has the
strongest regulatory orientation because it focuses on producer, manufacturer,
recycler, and refurbisher registration, EPR targets, certificates, returns, and
compliance monitoring. However, its public functions are not primarily designed
to coordinate household collection, transportation, repair, or user-facing
electronic waste classification.

Recykal provides broader circular economy and enterprise disposal services.
Its strengths include traceability, certified processing, compliance documentation,
recyclable-material transactions, and value recovery. Nevertheless, its
documented services primarily target organizations operating within an
established commercial recycling ecosystem and do not clearly provide the
complete set of role-specific workflows proposed for EcoTrace.

ecoATM demonstrates the practical value of automation and financial
incentives in encouraging device recovery. Its kiosks provide convenient device
assessment and immediate compensation, but the system is mainly limited to
eligible consumer devices and requires users to visit physical kiosk locations. It
does not provide a general-purpose electronic waste lifecycle platform for
households, businesses, collectors, drivers, technicians, recyclers, and
environmental authorities.

E-Tadweer demonstrates how mobile technology and rewards can increase
citizen participation in formal electronic waste recycling. Its drop-off-location and
voucher model is valuable for public engagement, but the reviewed information
does not describe complete reverse logistics, route optimization, repair,
refurbishment, or environmental monitoring workflows.

Green Grid provides the closest conceptual comparison to EcoTrace because it
combines pickup scheduling, mapping, rewards, educational content, AI
assistance, impact calculation, and marketplace functionality. However, it is
presented as a web-based academic prototype, and the publication does not
demonstrate the complete cross-platform and role-specific operational
architecture proposed for EcoTrace.

EcoTrace combines capabilities that are distributed across the reviewed
systems. Its design integrates electronic waste registration, AI-assisted image
classification, pickup scheduling, location and route services, reverse logistics,
repair, refurbishment, recycling, marketplace management, notifications,
reporting, and environmental monitoring within a single system.

The proposed system also supports Android, Web, and Windows applications
through Flutter, while the backend coordinates Cloud Firestore, Firebase
Authentication, Firebase Cloud Messaging, Cloudinary, Google Maps Platform,
Render, and PyTorch. This integrated architecture is intended to improve
accessibility, traceability, stakeholder coordination, and lifecycle visibility.

It is important to distinguish the proposed scope of EcoTrace from empirically
validated implementation results. The comparison in Table~\ref{tab:comparative-analysis}
reflects the approved system design and intended functional coverage.
Chapter~\ref{ch:implementation} documents which functions were fully
implemented, tested, partially implemented, or reserved for future development.

\section{Summary and Gaps in Literature}

\subsection{Summary of the Reviewed Literature}

This chapter reviewed the conceptual, technological, and practical foundations
relevant to the development of the EcoTrace system. The review established that
electronic waste is a rapidly increasing environmental and public health challenge
requiring coordinated collection, transportation, recovery, recycling, and
monitoring mechanisms.

The circular economy literature emphasizes extending product lifecycles
through reuse, repair, refurbishment, remanufacturing, and recycling rather than
relying on the traditional linear model of production, consumption, and disposal.
Reverse logistics complements this approach by coordinating the return of
end-of-life electronic products from consumers to collection centres, repair
facilities, refurbishment centres, recyclers, and manufacturers.

The reviewed literature also demonstrates that intelligent waste management
systems increasingly rely on cloud computing, artificial intelligence, geographic
information systems, mobile applications, data analytics, and real-time
communication. These technologies improve traceability, operational efficiency,
waste classification, collection planning, route optimization, stakeholder
coordination, and environmental decision-making.

The technology review examined Flutter, FastAPI, Cloud Firestore, Firebase
Authentication, Firebase Cloud Messaging, Cloudinary, Google Maps Platform,
Render, and PyTorch. Flutter supports the development of Android, Web, and
Windows applications from a single codebase. FastAPI provides high-performance
backend services, while Render supports cloud-based deployment of the backend
API. Cloud Firestore provides scalable document-based data storage and
real-time synchronization. Firebase Authentication supports secure user identity
management, while Firebase Cloud Messaging enables real-time notifications.
Cloudinary provides managed image storage and delivery, Google Maps Platform
supports location and route services, and PyTorch provides the artificial
intelligence framework required for electronic waste image classification.

The review of existing systems showed that the CPCB E-Waste EPR Portal,
Recykal, ecoATM, E-Tadweer, and Green Grid each address significant aspects of
electronic waste management. CPCB primarily emphasizes regulation and
compliance, Recykal focuses on commercial circularity and traceability, ecoATM
provides automated device take-back, E-Tadweer promotes citizen participation
through incentives, and Green Grid combines pickup services, mapping, rewards,
AI assistance, and marketplace functionality
\citep{cpcb2026,ecoatm2026a,etadweer2026,jagtap2026,recykal2026b}.

However, the reviewed systems differ considerably in purpose, target users,
geographical context, platform coverage, and lifecycle functionality. The publicly
documented evidence does not demonstrate that any single reviewed platform
integrates all the operational, technical, environmental, and stakeholder-management
capabilities proposed for EcoTrace.

\subsection{Identified Gaps in Literature and Existing Systems}

Based on the reviewed literature and systems, the following gaps were identified.

\textbf{Fragmented Electronic Waste Lifecycle Management.} Many reviewed
systems focus on a specific stage of electronic waste management, such as
regulatory compliance, device collection, recycling, or consumer incentives. The
public documentation generally does not demonstrate a unified platform covering
electronic waste registration, collection, transportation, repair, refurbishment,
recycling, marketplace operations, and environmental reporting as connected
lifecycle processes. EcoTrace addresses this gap by integrating the major
electronic waste recovery activities within a single information system.

\textbf{Limited Cross-Platform Accessibility.} Some existing solutions are
primarily web-based, mobile-oriented, or dependent on physical kiosks. This limits
accessibility for stakeholders whose operational requirements differ across
mobile, browser, and desktop environments. EcoTrace addresses this limitation
through Flutter applications designed for Android, Web, and Windows platforms
using a shared codebase.

\textbf{Limited AI-Based Waste Image Classification.} Although some
reviewed systems incorporate automated assessment or AI assistance, the
available documentation does not consistently demonstrate image-based
classification of multiple electronic waste categories as an integrated user-facing
service. EcoTrace proposed a PyTorch-based image classification component that
would analyse uploaded electronic waste images and return a predicted category
and confidence score. The final model architecture and measured performance
were to be documented after implementation and evaluation.

\textbf{Inadequate Role-Specific Operational Workflows.} The reviewed
systems generally target selected stakeholders such as producers, recyclers,
enterprises, consumers, or regulatory authorities. There is limited evidence of
integrated workflows designed specifically for households, businesses, collectors,
drivers, technicians, recyclers, environmental officers, administrators, and super
administrators within one platform. EcoTrace addresses this gap through
role-based interfaces, permissions, dashboards, and operational workflows for its
different actor categories.

\textbf{Limited Integration of Repair and Refurbishment.} Recycling receives
considerable attention in electronic waste systems, but repair and refurbishment
are not always implemented as detailed operational workflows. This limits the
ability of systems to support product-life extension before devices are dismantled
for material recovery. EcoTrace includes repair diagnosis, repair status tracking,
refurbishment records, inspection, quality assessment, and reusable device
management to support circular economy principles.

\textbf{Limited Route Optimization and Logistics Coordination.} Several
systems provide collection-point information or basic mapping, but the reviewed
public sources do not consistently demonstrate integrated route optimization,
collector assignment, driver coordination, transportation tracking, and
processing-facility delivery. EcoTrace proposed integrating Google Maps Platform
with pickup requests, collection locations, route visualization, and route
optimization to support reverse logistics operations.

\textbf{Limited Environmental Monitoring and Analytics.} Existing solutions
may provide regulatory reports, impact calculators, or transaction records, but
there is limited evidence of unified environmental monitoring dashboards
combining collection activities, recycling statistics, recovered materials,
operational performance, and environmental indicators. EcoTrace provides
reporting and analytics functions for administrators and environmental officers to
support monitoring, planning, compliance, and decision-making.

\textbf{Insufficient Integration of Cloud Services.} The literature demonstrates
the individual value of cloud databases, authentication services, messaging
platforms, media hosting, mapping services, and artificial intelligence. However,
the reviewed electronic waste systems do not consistently demonstrate the same
unified cloud architecture proposed for EcoTrace. EcoTrace integrates a cloud
backend, Render, Cloud Firestore, Firebase Authentication, Firebase Cloud
Messaging, Cloudinary, and a mobile-money payment gateway within one
coordinated architecture -- the exact composition of which is documented, and
reconciled against this chapter's original technology proposal, in
Chapter~\ref{ch:implementation}.

\textbf{Limited Adaptation to the Sierra Leonean Context.} Most reviewed
electronic waste platforms were developed for countries with different regulatory
environments, digital infrastructure, commercial recycling ecosystems, and
institutional arrangements. Their features may therefore not directly address the
operational conditions and stakeholder requirements of Sierra Leone. EcoTrace is
designed with Sierra Leone as its primary application context while maintaining
an architecture that can be adapted to other regions -- most concretely realised in
its mobile-money payment integration built specifically against Sierra Leone's
Orange Money and Afrimoney networks (Section~\ref{sec:payment-module}, Chapter~\ref{ch:implementation}).

\textbf{Limited Evidence of Complete End-to-End Validation.} Some reviewed
systems are commercial services, government portals, or academic prototypes.
Their public sources provide different levels of technical detail, and several do not
present complete evidence covering implementation, integration testing,
performance evaluation, user acceptance testing, and deployment across all
proposed functions. The EcoTrace study documents the implemented modules,
test cases, results, and limitations in Chapter~\ref{ch:implementation}. Features
that remained incomplete, or that took a different technical path from what is
proposed in this chapter, are identified explicitly rather than reported as fully
operational.

\subsection{Justification for the Proposed EcoTrace System}

The identified gaps justify the development of EcoTrace as an integrated,
intelligent, and cross-platform electronic waste recovery and circular economy
management system. The proposed system combines electronic waste
registration, AI-assisted classification, pickup scheduling, reverse logistics,
repair, refurbishment, recycling, marketplace management, notifications,
environmental monitoring, and reporting within one platform.

The system also integrates mobile, web, and Windows applications with a
cloud backend. Cloud Firestore provides centralized data management, Firebase
Authentication manages user identity, Firebase Cloud Messaging supports
real-time notifications, and Cloudinary hosts electronic waste images. This
integrated approach is intended to improve electronic waste traceability,
stakeholder coordination, operational efficiency, resource recovery,
environmental monitoring, and circular economy participation. It therefore
provides the conceptual and technical basis for the system analysis and design
presented in Chapter~\ref{ch:analysis-design}.

\subsection{Chapter Summary}

This chapter reviewed the concepts, technologies, and existing systems relevant
to intelligent electronic waste management. It examined electronic waste,
circular economy principles, reverse logistics, intelligent waste management,
cloud computing, artificial intelligence, and geographic information systems.

The chapter also reviewed the technologies proposed for EcoTrace and
compared five related electronic waste management systems. The analysis
identified gaps relating to lifecycle integration, cross-platform accessibility, AI
classification, role-specific workflows, repair and refurbishment, route
optimization, environmental analytics, cloud integration, local adaptation, and
implementation validation.

The next chapter presents the methodology, analysis, requirements, and design
of the proposed EcoTrace system.
